% Dokumentacja projektu — Wizja komputerowa (AGH)
% Kompilacja:  pdflatex dokumentacja.tex   (dwukrotnie, dla spisu i odnośników)
% UWAGA: limit 5 stron (z obrazami, tabelami i bibliografią włącznie).
%        Przekroczenie limitu = 0 pkt. Marginesy i odstępy są celowo zwarte.
\documentclass[10pt,a4paper]{article}

\usepackage[T1]{fontenc}
\usepackage[utf8]{inputenc}
\usepackage[polish]{babel}
\usepackage[margin=2cm]{geometry}
\usepackage{graphicx}
\usepackage{booktabs}
\usepackage{microtype}
\usepackage{siunitx}
\usepackage{caption}
\usepackage{subcaption}
\usepackage[hidelinks]{hyperref}

\captionsetup{font=small,labelfont=bf}
\setlength{\parskip}{2pt}
\setlength{\parindent}{0pt}

% Ścieżka do rysunków eksportowanych z notatnika (viz.* zwraca obiekt Figure,
% który można zapisać przez fig.savefig("docs/figures/...")).
\graphicspath{{figures/}}

\title{\vspace{-1.5em}Klasyfikacja obrazów metodami klasycznej wizji komputerowej}
\author{Hubert Miklas \and \textit{[współautor]} \\[2pt]
\small Akademia Górniczo-Hutnicza --- Wizja komputerowa}
\date{\vspace{-1em}\small \today}

\begin{document}
\maketitle
\vspace{-1.5em}

% ============================================================ %
\section{Wprowadzenie i zbiór danych}
% ============================================================ %
Celem projektu jest klasyfikacja obrazów wyłącznie metodami klasycznej wizji
komputerowej i uczenia maszynowego (bez sieci neuronowych). Potok obejmuje:
wczytanie danych, preprocessing, ekstrakcję cech, klasyfikację oraz dobór
parametrów na zbiorze treningowo-walidacyjnym i ocenę na osobnym zbiorze
testowym.

\textbf{Zbiór danych.} % TODO: opisać WYBRANY zbiór (1-2 zdania) + skąd pochodzi.
Pracujemy na zbiorze \texttt{[nazwa]} zawierającym \num{[N]} obrazów RGB
w \num{[K]} klasach. Kod jest skonfigurowany tak, by jednym przełącznikiem
(\texttt{get\_config(...)}) wybrać jeden z trzech zbiorów:
(i) liście ziemniaka \emph{early blight} vs \emph{healthy} (\num{305} obrazów),
(ii) liście ziemniaka \emph{early} vs \emph{late blight} (\num{2000} obrazów),
(iii) owoce \emph{jabłka} vs \emph{pomidory} (\num{391}, z gotowym podziałem
train/test). % TODO: zostawić tylko ten faktycznie użyty; resztę usunąć.

\begin{figure}[h]
  \centering
  % TODO: viz.show_class_samples(...) -> fig.savefig("figures/samples.png")
  % \includegraphics[width=0.95\linewidth]{samples.png}
  \fbox{\parbox[c][2.2cm][c]{0.95\linewidth}{\centering\small
    [Rys.~--- przykładowe obrazy z każdej klasy (\texttt{viz.show\_class\_samples})]}}
  \caption{Reprezentatywne próbki z poszczególnych klas zbioru.}
  \label{fig:samples}
\end{figure}

% ============================================================ %
\section{Preprocessing}
% ============================================================ %
% TODO: opisać realnie zastosowane kroki i dobrane parametry; usunąć niewykorzystane.
Zastosowano następujący łańcuch przetwarzania (moduł \texttt{preprocessing}):
\begin{itemize}\itemsep1pt
  \item \textbf{Usunięcie artefaktów} --- filtracja medianowa/bilateralna
        ograniczająca szum i artefakty kompresji JPEG.
  \item \textbf{Poprawa jasności i koloru} --- CLAHE na kanale \emph{L}
        (przestrzeń LAB) oraz korekcja balansu bieli metodą \emph{gray-world}.
  \item \textbf{Segmentacja obiektu od tła} --- progowanie [Otsu / w przestrzeni
        HSV / GrabCut], a następnie operacje morfologiczne (otwarcie + domknięcie)
        i wybór największej składowej spójnej.
\end{itemize}
% TODO: wstawić 1 rysunek "przed/po" + maskę (viz.show_pair / show_mask_overlay).

% ============================================================ %
\section{Ekstrakcja cech}
% ============================================================ %
Wyznaczamy trzy komplementarne rodziny cech (moduł \texttt{features}); każda ma
czytelną listę nazw ułatwiającą analizę ważności cech.
\begin{itemize}\itemsep1pt
  \item \textbf{Tekstura (macierz GLCM).} Po kwantyzacji do \num{32} poziomów
        liczymy GLCM dla odległości $d\in\{1,3\}$ i kątów
        $\{0,45,90,135\}^\circ$, a z niej deskryptory Haralicka
        (\emph{contrast, dissimilarity, homogeneity, energy, correlation, ASM})
        oraz entropię --- łącznie \num{56} cech.
  \item \textbf{Cechy koloru.} Średnia, odchylenie i skośność dla kanałów RGB
        oraz HSV liczone na obszarze obiektu (\num{18} cech).
  \item \textbf{Deskryptory kształtu.} Z największego konturu: pole, obwód,
        zwartość, mimośród, \emph{extent}, średnica równoważna, bbox oraz
        statystyki promieni (\num{11} cech).
\end{itemize}
% TODO: które rodziny ostatecznie użyto (use=...) i dlaczego (np. rozkłady cech).
% TODO: wstawić viz.plot_feature_distributions dla 3-4 najbardziej różnicujących cech.

% ============================================================ %
\section{Klasyfikacja i dobór parametrów}
% ============================================================ %
Każdy model to potok \texttt{StandardScaler} $+$ klasyfikator (skalowanie uczone
tylko na foldzie treningowym, bez wycieku informacji). Porównujemy klasyfikatory
klasyczne: \textbf{SVM} i \textbf{drzewo decyzyjne} (wymagane w zadaniu) oraz
baseline'y: las losowy, kNN i regresję logistyczną.

Hiperparametry dobrano metodą \texttt{GridSearchCV} z walidacją krzyżową
(\emph{stratified $k$-fold}, metryka \texttt{f1\_macro}) na zbiorze
treningowo-walidacyjnym. % TODO: podać finalne siatki i najlepsze parametry.
Przeszukiwano m.in. $C$, $\gamma$ i jądro dla SVM oraz głębokość i
\texttt{min\_samples\_leaf} dla drzewa.

% ============================================================ %
\section{Wyniki i krytyczna analiza}
% ============================================================ %
Końcową ocenę przeprowadzono \emph{raz}, na wydzielonym zbiorze testowym.
Tabela~\ref{tab:wyniki} zestawia metryki; raportujemy średnie makro
(traktują klasy równo) oraz ważone.

\begin{table}[h]
  \centering\small
  \caption{Metryki na zbiorze testowym. % TODO: uzupełnić liczbami z metrics.compare_models
  }
  \label{tab:wyniki}
  \begin{tabular}{lcccc}
    \toprule
    Model & Accuracy & Precision\textsubscript{macro} & Recall\textsubscript{macro} & F1\textsubscript{macro} \\
    \midrule
    SVM              & -- & -- & -- & -- \\
    Drzewo decyzyjne & -- & -- & -- & -- \\
    Las losowy       & -- & -- & -- & -- \\
    \bottomrule
  \end{tabular}
\end{table}

\begin{figure}[h]
  \centering
  % TODO: viz.plot_confusion_matrix(...) -> fig.savefig("figures/confusion.png")
  \fbox{\parbox[c][2.6cm][c]{0.5\linewidth}{\centering\small
    [Rys.~--- macierz pomyłek najlepszego modelu]}}
  \caption{Macierz pomyłek najlepszego klasyfikatora na zbiorze testowym.}
  \label{fig:cm}
\end{figure}

\textbf{Analiza.} % TODO: napisać autorskie wnioski (najważniejsza część oceny):
% - która rodzina cech daje największy wkład (feature importance / rozkłady)?
% - które klasy są mylone i dlaczego (off-diagonal macierzy pomyłek)?
% - rozbieżność accuracy vs F1_macro przy niezbalansowanych klasach?
% - wpływ preprocessingu/segmentacji na wynik (ablacja)?
[do uzupełnienia po eksperymentach]

% ============================================================ %
\section{Wnioski}
% ============================================================ %
% TODO: 3-5 zdań podsumowania: najlepsza konfiguracja, ograniczenia, co dalej.
[do uzupełnienia]

% ============================================================ %
\begin{thebibliography}{9}
\small
\bibitem{haralick} R.~M. Haralick, K.~Shanmugam, I.~Dinstein,
  \emph{Textural Features for Image Classification}, IEEE TSMC, 1973.
\bibitem{skimage} S.~van der Walt i in.,
  \emph{scikit-image: image processing in Python}, PeerJ, 2014.
\bibitem{sklearn} F.~Pedregosa i in.,
  \emph{Scikit-learn: Machine Learning in Python}, JMLR, 2011.
\bibitem{opencv} G.~Bradski, \emph{The OpenCV Library}, Dr.~Dobb's Journal, 2000.
% TODO: dodać źródło zbioru danych (np. PlantVillage) oraz pozostałe odnośniki.
\end{thebibliography}

\end{document}
